% Options for packages loaded elsewhere
\PassOptionsToPackage{unicode}{hyperref}
\PassOptionsToPackage{hyphens}{url}
\documentclass[
  11pt,
]{article}
\usepackage{xcolor}
\usepackage[margin=25mm]{geometry}
\usepackage{amsmath,amssymb}
\setcounter{secnumdepth}{-\maxdimen} % remove section numbering
\usepackage{iftex}
\ifPDFTeX
  \usepackage[T1]{fontenc}
  \usepackage[utf8]{inputenc}
  \usepackage{textcomp} % provide euro and other symbols
\else % if luatex or xetex
  \usepackage{unicode-math} % this also loads fontspec
  \defaultfontfeatures{Scale=MatchLowercase}
  \defaultfontfeatures[\rmfamily]{Ligatures=TeX,Scale=1}
\fi
\usepackage{lmodern}
\ifPDFTeX\else
  % xetex/luatex font selection
\fi
% Use upquote if available, for straight quotes in verbatim environments
\IfFileExists{upquote.sty}{\usepackage{upquote}}{}
\IfFileExists{microtype.sty}{% use microtype if available
  \usepackage[]{microtype}
  \UseMicrotypeSet[protrusion]{basicmath} % disable protrusion for tt fonts
}{}
\makeatletter
\@ifundefined{KOMAClassName}{% if non-KOMA class
  \IfFileExists{parskip.sty}{%
    \usepackage{parskip}
  }{% else
    \setlength{\parindent}{0pt}
    \setlength{\parskip}{6pt plus 2pt minus 1pt}}
}{% if KOMA class
  \KOMAoptions{parskip=half}}
\makeatother
\setlength{\emergencystretch}{3em} % prevent overfull lines
\providecommand{\tightlist}{%
  \setlength{\itemsep}{0pt}\setlength{\parskip}{0pt}}
\usepackage{bookmark}
\IfFileExists{xurl.sty}{\usepackage{xurl}}{} % add URL line breaks if available
\urlstyle{same}
\hypersetup{
  hidelinks,
  pdfcreator={LaTeX via pandoc}}

\author{}
\date{}

\begin{document}

\section{Deformations and the analytic index in
E-theory}\label{deformations-and-the-analytic-index-in-e-theory}

\emph{Written by GPT-6.1 Sol (OpenAI) in Codex, Ultra setting, October
2026. Public domain (CC0).}

A deformation replaces an operator algebra by its limiting symbol
algebra. Evaluation at a small positive parameter almost preserves
multiplication on lifts of symbols. Those evaluations define an
asymptotic morphism. For the tangent groupoid, the resulting map on
K-theory is the analytic index.

Algebras are complex and separable. Write \(SA=C_0((0,1),A)\),
\(\mathcal K=\mathcal K(\ell^2)\), and
\(E(A,B)=[[SA,SB\otimes\mathcal K]]\), with products in the order of
composition. The definition of asymptotic morphism, composition,
homotopy and stabilization are proved in
\href{KT-KK-21.html\#1-families-quotients-and-continuous-representatives}{\emph{Asymptotic
morphisms and E-theory}, Sections 1--4}. Its
\href{KT-KK-21.html\#6-transporting-the-bott-equivalence}{Lemma 6.1 and
equation (6.6)} prove the product-preserving functor
\(\mathfrak c:KK\to E\). The deformation and K-theory arguments below do
not require contravariant E-theory exactness.

\subsection{1. A deformation and its continuous
lifts}\label{a-deformation-and-its-continuous-lifts}

A \textbf{deformation extension} is an exact sequence \[
0\longrightarrow I=C_0((0,1],B)\longrightarrow D
 \xrightarrow{q}A\longrightarrow0,
\tag{22.1}
\] together with positive evaluations \(e_\tau:D\to B\),
\(0<\tau\leq1\), whose restriction to \(I\) is ordinary evaluation. We
require \(\tau\mapsto e_\tau(d)\) to be norm continuous at positive
parameters, and the evaluations together with \(q\) to be faithful. A
continuous \(C([0,1])\)-algebra with zero fibre \(A\) and trivial
positive field \(B\) has these properties. The proof also works for an
upper semicontinuous field satisfying (22.1).

For \(i\in I\), \[
\|e_\tau(i)\|\longrightarrow0\quad(\tau\downarrow0).
\tag{22.2}
\] The convergence is uniform on compact subsets of \(I\): cover such a
set by finitely many norm balls of radius \(\epsilon\), apply (22.2) to
their centres and use contractivity of evaluation. This elementary
observation will control homotopies as well as individual errors.

\textbf{Lemma 22.1 (a continuous lift).} There is a continuous,
generally nonlinear map \(\sigma:A\to D\), with \(q\sigma(a)=a\),
\(\sigma(0)=0\) and \(\|\sigma(a)\|\leq4\|a\|\).

\textbf{Proof.} The quotient norm gives a lift of every nonzero \(a\) of
norm at most \(2\|a\|\). Choose a countable dense family \(a_j\) in the
unit sphere and such lifts \(d_j\) with \(\|d_j\|\leq2\). For a unit
vector \(a\), put \[
w_j(a)=2^{-j}\max\{0,\tfrac12-\|a-a_j\|\},\qquad
H(a)=\frac{\sum_jw_j(a)d_j}{\sum_jw_j(a)}.
\tag{22.3}
\] The denominator is positive. It is bounded below on a neighbourhood
of each unit vector; the two series converge uniformly there. Thus \(H\)
is continuous, has norm at most two, and \(\|qH(a)-a\|\leq1/2\). Extend
\(H\) by positive radial homogeneity and set \(H(0)=0\). Starting with
\(r_0(a)=a\), set \(r_{k+1}(a)=r_k(a)-qH(r_k(a))\). The residual norm is
at most \(2^{-k}\|a\|\). Hence \[
\sigma(a)=\sum_{k=0}^{\infty}H(r_k(a))
\tag{22.4}
\] converges uniformly on bounded sets, is continuous, has norm at most
\(4\|a\|\), and telescoping gives \(q\sigma(a)=a\). If \(A=0\), use the
zero map. \(\square\)

\textbf{Theorem 22.2 (the deformation class).} The family \[
\alpha_t(a)=e_{1/t}(\sigma(a)),\qquad t\geq1,
\tag{22.5}
\] is an asymptotic morphism \(A \rightsquigarrow B\). Its asymptotic
equivalence class is independent of \(\sigma\), and its suspension
followed by a rank-one corner defines a canonical \(d_D\in E(A,B)\).

\textbf{Proof.} The maps are bounded at each \(a\), and continuous in
\(t\). Every linearity, adjoint or product error of \(\sigma\) is in
\(\ker q=I\). For instance \[
\alpha_t(ab)-\alpha_t(a)\alpha_t(b)
=e_{1/t}\bigl(\sigma(ab)-\sigma(a)\sigma(b)\bigr)\longrightarrow0.
\tag{22.6}
\] The other two identities have the same proof. Two lifts differ at
each \(a\) by an element of \(I\), so their evaluated families are
asymptotically equivalent. Their linear interpolation is an asymptotic
homotopy: its errors are finite polynomials in the interpolation
parameter with coefficients in \(I\), giving uniform convergence by
(22.2).

Here is a representative on the whole suspension, rather than just on
elementary tensors. For \(f\in SA\) set
\((S\sigma)(f)(u)=\sigma(f(u))\). Continuity of \(\sigma\) and its norm
bound make this an element of \(SD\). Each of its algebraic defects is
in \(C_0((0,1),I)\). Uniform approximation by finitely many compactly
supported \(I\)-valued functions, and (22.2), show that applying
\(Se_{1/t}\) makes these defects vanish in the \(SB\) norm.
Positive-time continuity holds uniformly in \(u\), because the compact
image of \(f\) under \(\sigma\) can be covered by finite norm nets.
Compose with \(b\mapsto b\otimes e_{11}\). This gives an asymptotic
morphism \(SA \rightsquigarrow SB\otimes\mathcal K\). The same
interpolation proves independence on the suspension. The different
rank-one corners are related by a unitary path on their
finite-dimensional span; conjugating by that path gives the same stable
class. \(\square\)

Only errors lying in the cone ideal were used. In particular, this proof
does not assert a completely positive lift of an arbitrary quotient.

\subsection{2. The map on K-theory}\label{the-map-on-k-theory}

We prove explicitly why the zero evaluation in (22.1) is a
K-isomorphism. A positive-time field may fail to converge in the fixed
algebra \(B\) as \(\tau\downarrow0\); the argument will never require
that convergence.

\textbf{Lemma 22.3 (freezing the positive field).} For \(0<\delta\leq1\)
there is a homomorphism \(R_\delta:D\to D\) satisfying \[
qR_\delta=q,\qquad
e_\tau R_\delta(d)=e_{\min(\tau,\delta)}(d).
\tag{22.7}
\] It is homotopic to the identity.

\textbf{Proof.} The difference \(e_{\min(\tau,\delta)}(d)-e_\tau(d)\) is
a continuous \(B\)-valued function which is identically zero for
\(\tau\leq\delta\); it is therefore an element of \(I\). Add it to \(d\)
to define \(R_\delta(d)\). The displayed evaluations show that it
preserves addition, products and adjoints. Their joint faithfulness
proves those identities in \(D\). The same faithfulness gives its norm
bound by the supremum of the original fibre norms. Allowing \(\delta\)
to move from its prescribed value to one is point-norm continuous: below
its smallest value the difference is zero, and on the remaining compact
positive interval uniform continuity of the original field controls the
supremum norm. At one the map is the identity. \(\square\)

\textbf{Proposition 22.4.} The map \(q_*:K_j(D)\to K_j(A)\) is an
isomorphism, \(j=0,1\), and the K-map of (22.5) is \[
(d_D)_*=(e_1)_*q_*^{-1}.
\tag{22.8}
\]

\textbf{Proof in degree zero.} Apply the lift entrywise to matrices and
extend it to unitizations by leaving scalar matrices fixed. If
\(p\in M_N(A^+)\) is a projection, a self-adjoint lift \(r\in M_N(D^+)\)
has \(r^2-r\in M_N(I)\). Thus at all sufficiently small positive
parameters its spectrum has two disjoint clusters near zero and one.
Choose a continuous real function \(h\) which is zero near the first
cluster and one near the second. Then \(qh(r)=p\), and \(e_\tau h(r)\)
is a projection for \(0<\tau\leq\delta\), after making \(\delta\)
smaller if necessary. Its scalar part is a projection as well. Lemma
22.3 now makes \(R_\delta(h(r))\) an actual projection lifting \(p\).
This proves surjectivity on the stabilized projection group.

For injectivity suppose two stabilized projections in \(D^+\) have the
same K-class after applying \(q\). By the defining stabilized projection
relations, add a common projection and, if needed, scalar projections to
obtain a continuous projection path between their images. Lift the added
projection as above. Lift the path continuously using Lemma 22.1,
self-adjointize it and use the same function \(h\). Its error path is a
compact subset of \(M_N(I)\), so a single \(\delta\) works for the whole
path. Applying \(R_\delta\) gives a projection path in \(D^+\). At its
endpoints this lifted path differs from the frozen original projections
by a norm tending uniformly to zero as \(\delta\downarrow0\): their
differences have quotient zero, and near zero the fixed functional
calculus preserves the projection clusters. Choose the norm difference
less than \(1/2\). The spectral projection of the straight interpolation
gives a projection homotopy between each endpoint pair. Finally
\(R_\delta\) is homotopic to the identity. Thus the original K-classes
are equal. This also respects the scalar-part map, so restriction from
unitizations gives injectivity for \(K_0(D)\).

For clarity, the K-map of the asymptotic family has the usual direct
representative. For a projection \(p\), self-adjointize its matrix image
under \(\alpha_t\) and take the spectral projection above \(1/2\) for
sufficiently large \(t\). The adjoint and product errors tend to zero,
so the gap is present. A stabilized projection path gives a uniform gap
by the compact-set argument, proving that this construction depends only
on its K-class. If \(P\) is an actual lift of \(p\), the matrix image of
\(p\) differs from \(e_{1/t}(P)\) by a matrix in \(I\), evaluated near
zero. Their corrected projections are consequently homotopic. The
continuous positive path \(e_\tau(P)\), \(1/t\leq\tau\leq1\), identifies
their K-class with \((e_1)_*[P]\). This proves (22.8).

This is also the action of the E-class under \(E(\mathbb C,A)=K_0(A)\).
In the positive suspension convention a projection gives the unitary
loop \(1+(e^{2\pi iu}-1)p\). Composing its suspended representative with
our family gives, up to errors uniform in \(u\), the loop
\(1+(e^{2\pi iu}-1)\alpha_t(p)\): asymptotic scalar linearity is uniform
on this compact scalar interval. Replacing the almost projection by its
spectral correction changes the loop by a norm tending to zero. Polar
correction and the logarithmic path between nearby unitaries leave its
K-class unchanged. It is therefore the positive loop of precisely the
projection class calculated above.

\textbf{Degree one.} Lift a unitary in a matrix unitization. At small
positive parameters both \(v^*v-1\) and \(vv^*-1\) are small. Multiply
the lift \(v\) by \(g(v^*v)\), where \(g\) is a bounded continuous
function agreeing with \(x^{-1/2}\) near one. This expression is unitary
at the small fibres and at zero. Freezing makes it unitary on the whole
field. A compact unitary path has one common small interval, exactly as
above. Nearby unitaries are homotopic by polar normalization of their
straight interpolation; freezing is homotopic to the identity. These
facts prove bijectivity in degree one and identify the asymptotic K-map
with (22.8). \(\square\)

The proofs also show naturality under homomorphisms of deformation
extensions: quotient, positive evaluations and freezing commute with
such homomorphisms, and two choices of lifts have vanishing differences.

\textbf{Lemma 22.5 (a KK comparison when available).} If (22.1) is
semisplit, its zero evaluation is a KK-equivalence by the separable
semisplit exactness proved in
\href{KT-KK-14.html\#4-excision-and-the-six-term-sequences}{\emph{Exact
sequences in KK and the universal coefficient theorem}, Theorems
4.1--4.2}. In this case \[
d_D=\mathfrak c\bigl([q]^{-1}\otimes_D[e_1]\bigr).
\tag{22.9}
\]

\textbf{Proof.} The cone \(I\) contracts by \(i(\tau)\mapsto i(s\tau)\),
with value zero at \(s=0\); uniform continuity of the zero extension
makes this a point-norm homotopy. Semisplit KK exactness therefore makes
postcomposition by \([q]\) an isomorphism. Surjectivity supplies a right
inverse in \(KK(A,D)\), and injectivity with source \(D\) makes it a
left inverse as well.

The direct deformation family composed with \(q\) differs from
\(e_{1/t}\) at each \(d\) by \(e_{1/t}(\sigma(qd)-d)\), which tends to
zero. The exact homomorphisms \[
H_t(d)(s)=e_{s+(1-s)/t}(d),\qquad 0\leq s\leq1,
\tag{22.10}
\] give an asymptotic homotopy from \(e_{1/t}\) to \(e_1\). Every
algebraic defect is zero, and positive-interval uniform continuity gives
the required continuity in \(t\) and \(s\). Suspend and stabilize this
homotopy. Hence \(\mathfrak c([q])\otimes_A d_D=\mathfrak c([e_1])\).
The product-preserving functor takes the actual two KK inverse
identities to two E inverse identities, proving (22.9). No E-theory
half-exactness or cancellation theorem is used. \(\square\)

\subsection{3. The tangent groupoid and the analytic
index}\label{the-tangent-groupoid-and-the-analytic-index}

Let \(M\) be a Hausdorff second-countable smooth manifold without
boundary, with a positive density. Its tangent groupoid is \[
\mathbb T M=(TM\times\{0\})\ \sqcup\
(M\times M\times(0,1]).
\tag{22.11}
\] At zero its arrows compose by addition in each tangent space; at
positive parameter they compose as pairs. In a coordinate chart its
boundary coordinate is \((x,v,\tau)\mapsto(x,x-\tau v,\tau)\). The
divided-difference proof of chart compatibility, Hausdorffness and
smooth groupoid operations is \href{../KT-CP/KT-CP-17.html}{\emph{The
tangent groupoid and deformation to the normal cone}, Lemma 17.1 and
Proposition 17.2}. The Haar system is the linear density at zero and
\(\tau^{-n}d\mu\) at positive parameters, as proved in its Lemma 17.3.

Put \(D_M=C^*(\mathbb T M)\). The full invariant-restriction theorem and
finite-rank pair-kernel proof are
\href{../KT-CP/KT-CP-14.html}{\emph{Groupoid C}-algebras: full and
reduced*, Theorem 14.5 and Proposition 14.6}. Fibrewise Fourier
transformation at zero and that pair calculation give \[
0\longrightarrow C_0((0,1],\mathcal K(L^2M))
\longrightarrow D_M\xrightarrow{e_0}C_0(T^*M)
\longrightarrow0.
\tag{22.12}
\] For a nonempty manifold the pair Hilbert space is nonzero; an empty
manifold gives the zero class throughout. Positive trivialization
includes the factor \(\tau^{-n}\) prescribed by the Haar density. These
are precisely the algebra and fibre identifications in Proposition 17.4
of the tangent-groupoid lesson. Its Theorem 17.6 proves norm continuity
at zero: an upper bound comes from the quotient field, and the lower
bound uses localized Gaussian packets of spatial width \(\sqrt{\tau}\)
and phase \(e^{i\eta x/\tau}\). That proof also applies on noncompact
\(M\), since each packet lies in one relatively compact chart.

Theorem 22.2 therefore gives \[
\operatorname{Ind}^{E}_M\in
E(C_0(T^*M),\mathcal K(L^2M)).
\tag{22.13}
\] Use a rank-one Hilbert-space Morita corner when a scalar target is
desired. Proposition 22.4 identifies its K-map with \[
\operatorname{ind}_M=(e_1)_*(e_0)_*^{-1}:
K_0(C_0(T^*M))\longrightarrow K_0(\mathcal K(L^2M))
\cong\mathbb Z.
\tag{22.14}
\] This is the deformation analytic index of Theorem 17.7 in the
tangent-groupoid lesson.

We verify its operator meaning for closed manifolds using the complete
free-source quantization proof in
\href{KT-KK-15.html\#9-a-universal-quantization-morphism}{\emph{Dirac
classes and the cotangent Dolbeault element}, Lemma CI.10 and Theorems
CI.11--CI.14}. This comparison uses the classical order-zero index
statement proved there.

\textbf{Proposition 22.6 (the closed-manifold comparison).} For closed
\(M\), the semiclassical algebra \(\mathcal A_M\) of Lemma CI.10 is
\(D_M\), with the same two evaluations. Thus the E-class is the image of
its KK quantization morphism. For every classical elliptic order-zero
\(P\), \[
(\operatorname{Ind}^E_M)_*[\sigma_P]
=\operatorname{index}P.
\tag{22.15}
\]

\textbf{Proof.} Work first in a protected chart with half-densities. For
\(a\in C_c^\infty(T^*M)\), its semiclassical kernel is the inverse
Fourier transform of \(a(x,\xi)\) in the variable \(v=(x-y)/\tau\), with
the chart cutoffs of CI.10. It is smooth and rapidly decreasing in
\(v\), uniformly on compact base sets and \(0\leq\tau\leq1\). Cutting it
off at large \(|v|\) converges in the groupoid \(I\)-norm: both the row
and column integrals of the rapidly decreasing tail tend uniformly to
zero. The zero Fourier value is \(a\). Thus every quantization
generator, together with the positive cone ideal, belongs to \(D_M\).

Conversely a smooth compact groupoid kernel in a boundary chart has a
smooth zero kernel of compact \(v\)-support. Its Fourier transform
\(a_0(x,\xi)\) is rapidly decreasing in \(\xi\). Approximate it by
compact-frequency smooth symbols; Fourier integration by parts bounds
the inverse-transform error in the two \(I\)-integrals by finitely many
weighted derivative errors, which tend to zero under those cutoffs. The
original positive kernel differs near zero from the chart quantization
of \(a_0\) by norm \(O(\tau)\): its smooth parameter dependence, the
half-density factor and the source-coordinate cutoff each have a first
Taylor error bounded by \(\tau\) times an integrable compact or rapidly
decreasing \(v\)-majorant. The finite partition from CI.10 handles all
boundary charts. Kernels supported away from zero belong to the already
included compact cone ideal. Smooth compact kernels are dense in the
groupoid completion by chart convolution and compact cutoffs. Therefore
the two closed algebras coincide, with identical evaluations.

CI.10's coherent-vector section proves semisplitting. Lemma 22.5 now
identifies the E-class with the image of \([e_0]^{-1}[e_1]\). CI.11
identifies the latter, followed by Hilbert-space Morita, with the
\(J_-\) cotangent Dolbeault element and proves the full-source symbol
identity. Its scalar pairing is exactly \(\operatorname{index}P\), by
CI.9 and CI.14. Proposition 22.4 gives the same pairing directly in
K-theory, proving (22.15). \(\square\)

For a noncompact manifold (22.14) remains the deformation index on
compactly supported symbol K-classes. A bare elliptic operator on a
noncompact manifold need not be Fredholm; the formula does not assign a
Fredholm integer to such an operator without additional control at
infinity.

\subsection{4. A Euclidean calculation with its bivariant
conclusion}\label{a-euclidean-calculation-with-its-bivariant-conclusion}

For \(M=\mathbb R^n\), the algebra in (22.12) has zero fibre
\(C_0(\mathbb R^{2n})\). The inverse Bott assertion requires a bivariant
proof. Its K-index alone will not be used as an E inverse criterion.

\textbf{Lemma 22.7 (the positive Euclidean section).} The Euclidean zero
evaluation is semisplit.

\textbf{Proof.} Let \[
g_{z,\eta,\tau}(x)=(\pi\tau)^{-n/4}
\exp[-|x-z|^2/(2\tau)]\exp[i\eta\cdot(x-z)/\tau].
\tag{22.16}
\] Fourier unitarity in \(\eta\), followed by the normalized Gaussian
integral in \(z\), gives \[
\int |g_{z,\eta,\tau}\rangle\langle g_{z,\eta,\tau}|
\frac{dz\,d\eta}{(2\pi\tau)^n}=1.
\tag{22.17}
\] Consequently \(S_\tau(a)=\int a(z,\eta)|g_{z,\eta,\tau}\rangle
\langle g_{z,\eta,\tau}|\,dz\,d\eta/(2\pi\tau)^n\) is completely
positive and contractive. For a compactly supported smooth \(a\),
Gaussian convolution and Taylor's integral formula compare this
anti-Wick kernel to its semiclassical Fourier kernel with norm error
\(O(\sqrt{\tau})\). Indeed the first spatial and frequency increments
have Gaussian first moments \(O(\sqrt{\tau})\); the finite derivative
Schur bound from Lemma CI.10 controls the remaining kernel. Both
quantizations are compact at positive parameter: the Fourier kernel is
square-integrable, while the positive kernel is a finite integral of
rank-one operators on each compact phase support. The Fourier-kernel and
\(I\)-norm approximation proof in Proposition 22.6 applies here to
compact base supports, and every compact positive cone field is again in
\(D_{\mathbb R^n}\). Therefore \((a,S_\tau(a))\) is a section of that
algebra with zero value \(a\). Contractivity extends it to every
\(a\in C_0(\mathbb R^{2n})\), and positivity at every matrix level
survives norm closure. This is the required completely positive section.
\(\square\)

On \(L^2(\mathbb R^n,\Lambda\mathbb C^n)\), with exterior degree
grading, write \(\varepsilon_j\) and \(\iota_j\) for creation and
contraction. The operator \[
\mathcal Q_\tau=\sum_j
\bigl((x_j+\tau\partial_j)\varepsilon_j
 +(x_j-\tau\partial_j)\iota_j\bigr),\qquad \tau>0,
\tag{22.18}
\] has \[
\mathcal Q_\tau^2=-\tau^2\Delta+|x|^2
 +\tau(2N-n),\qquad N=\sum_j\varepsilon_j\iota_j.
\tag{22.19}
\] Its maximal weighted domain, essential self-adjointness, Hermite
completeness and spectral gap are the explicit oscillator proof of
\href{KT-KK-12.html\#4-the-first-product-and-its-gaussian}{\emph{Bott
periodicity in KK: the Bott and Dirac elements}, Lemma 4.0 and the first
Gaussian product}. The same proof after dilation applies for each
\(\tau>0\). At \(\tau=1\), its only kernel is the even normalized scalar
Gaussian; the next eigenvalue of its square is at least two.

We verify that its graph projection is a section, including the
unbounded spatial potential. This avoids an implicit compactification
theorem for \(\mathbb R^n\).

\textbf{Lemma 22.8 (heat kernels in the deformation algebra).} For each
\(s>0\), the family \(e^{-s\mathcal Q_\tau^2}\) is a section of
\(D_{\mathbb R^n}\) with zero value \(e^{-s(|x|^2+|\eta|^2)}\). The
family \(\mathcal Q_\tau e^{-s\mathcal Q_\tau^2}\) is also a section,
with zero value \((c(x)+f(\eta))e^{-s(|x|^2+|\eta|^2)}\), where
\(c=\varepsilon+\iota\) and \(f=i(\varepsilon-\iota)\).

\textbf{Proof.} The scalar kernel for \(-\tau^2\Delta+|x|^2\) is \[
K_{s,\tau}(x,y)=
(2\pi\tau\sinh(2\tau s))^{-n/2}
\exp\!\left[-\frac{(|x|^2+|y|^2)\cosh(2\tau s)-2x\cdot y}
 {2\tau\sinh(2\tau s)}\right].
\tag{22.20}
\] Here is a direct verification. Write its exponent as
\(-a(s)(|x|^2+|y|^2)+b(s)x\cdot y\), where
\(a=(2\tau)^{-1}\coth(2\tau s)\) and \(b=(\tau\sinh(2\tau s))^{-1}\).
Differentiation gives \(a'=1-4\tau^2a^2=-\tau^2b^2\), \(b'=-4\tau^2ab\),
and the logarithmic derivative of its prefactor is \(-2n\tau^2a\).
Substitution into \(\partial_sK=(\tau^2\Delta_x-|x|^2)K\) proves that
equation. Gaussian integration gives its row mass
\((\cosh(2\tau s))^{-n/2}\exp[-\tanh(2\tau s)|x|^2/(2\tau)]\leq1\);
symmetry gives the column bound. The elementary Schur estimate, obtained
by Cauchy--Schwarz against each row measure and then integration in the
other variable, makes the kernel operators contractions. As
\(s\downarrow0\), the Gaussian mass approaches one near \(x=y\) and its
tails vanish, so they converge to the identity on compactly supported
smooth vectors. The family with this kernel equals the spectral heat
semigroup: for the difference of the two solutions, the derivative of
its squared Hilbert norm is \(-2\langle Hu,u\rangle\leq0\); start at a
positive small time and let that time tend to zero. Smooth Gaussian
truncations justify the maximal-domain calculation. Density and the
contraction bound extend the identity to \(L^2\). Formula (22.19)
multiplies the scalar kernel by the finite matrix \(e^{-s\tau(2N-n)}\).

In tangent coordinates \(y=x-\tau v\), the Haar-normalized kernel is
\(\tau^nK_{s,\tau}(x,x-\tau v)e^{-s\tau(2N-n)}\). Its scalar part is \[
\left(\frac{\tau}{2\pi\sinh(2\tau s)}\right)^{n/2}
\exp\!\left[-\frac{\tanh(\tau s)}{\tau}
 |x-\tfrac{\tau v}{2}|^2-\frac{\tau\coth(\tau s)}4|v|^2\right].
\tag{22.21}
\] The ratios in this formula extend smoothly to \(\tau=0\). At zero it
becomes \[
(4\pi s)^{-n/2}\exp[-s|x|^2-|v|^2/(4s)],
\tag{22.22}
\] whose fibrewise Fourier transform is the asserted scalar symbol.

For fixed \(s>0\), the two positive quadratic coefficients in (22.21)
have positive lower bounds on \(0\leq\tau\leq1\). The change
\((x,v)\mapsto(x-\tau v/2,v)\) and its inverse have uniformly bounded
matrices there. Consequently the kernel, and each of the finitely many
derivatives used below, are bounded by a polynomial times
\(C_s e^{-c_s(|x|^2+|v|^2)}\), with \(c_s>0\). Cutting off both \(x\)
and \(v\) therefore converges uniformly in the row and column \(I\)-norm
integrals; inversion changes \(x\) to \(x-\tau v\) and leaves the same
type of bound. The cut-off functions are compact groupoid kernels. Thus
(22.21) defines a section of its C*-algebra.

On every compact \(s\)-interval inside \((0,\infty)\), these majorants
and their first \(s\)-derivatives are uniform. Dominated convergence in
the two \(I\)-integrals therefore proves norm continuity in \(s\), as
required for the integrals below.

Apply \(\mathcal Q_\tau\) to the left kernel variable, holding \(y\)
fixed. In \((x,v)\) coordinates the derivative \(\tau\partial_x\) is
\(\tau\partial_x+\partial_v\). Its application to (22.21) is another
polynomial Gaussian with the same uniform cutoff bounds. At zero Fourier
transformation takes \(\partial_v\) to \(i\eta\); hence the zero symbol
is \(c(x)+f(\eta)\) times the scalar heat symbol. This proves the second
assertion, with its sign. \(\square\)

\textbf{Theorem 22.9 (the Euclidean Bott inverse).} With the outward
Bott position symbol \(c(x)+f(\eta)\), the E-theory index of
\(\mathbb R^n\), followed by Morita equivalence, is the inverse of that
Bott element.

\textbf{Proof.} Spectral calculus gives uniformly in \(0\leq\tau\leq1\)
\[
\|e^{-s\mathcal Q_\tau^2}\|\leq1,\qquad
\|\mathcal Q_\tau e^{-s\mathcal Q_\tau^2}\|
\leq(2es)^{-1/2}.
\] Thus the norm-convergent integrals \[
R=\int_0^\infty e^{-s}e^{-s\mathcal Q^2}\,ds,\qquad
V=\int_0^\infty e^{-s}\mathcal Q e^{-s\mathcal Q^2}\,ds
\tag{22.23}
\] are sections in the matrix algebra over \(D_{\mathbb R^n}\). At
positive parameters they are \((1+\mathcal Q_\tau^2)^{-1}\) and
\(\mathcal Q_\tau(1+\mathcal Q_\tau^2)^{-1}\); the zero values are the
same rational functions of \(c(x)+f(\eta)\).

Let \(A_\tau\) denote the part of \(\mathcal Q_\tau\) from even to odd
forms, and let \(R_+\), \(R_-\) be the two diagonal blocks of \(R\). The
projection \[
P_\tau=
\begin{pmatrix}
R_{+,\tau}&(A_\tau R_{+,\tau})^*\\
A_\tau R_{+,\tau}&1-R_{-,\tau}
\end{pmatrix},
\qquad P_\infty=\begin{pmatrix}0&0\\0&1\end{pmatrix},
\tag{22.24}
\] is the graph projection. Its difference from \(P_\infty\) belongs to
a matrix algebra over \(D_{\mathbb R^n}\), by (22.23). The
inverse-resolvent identities prove \(P_\tau^2=P_\tau=P_\tau^*\) at every
fibre, and joint faithfulness proves the identities in the algebra. At
zero this is exactly the graph class of the outward Bott symbol with
coordinate \(x+i\eta\), by the relative multiplication-cycle
construction in Theorems RK.3--RK.4 of the earlier difference-bundle
lesson. Call that class \(\beta_n\).

At one, replace \(A_1\) by \(\lambda A_1\), \(1\leq\lambda<\infty\). On
its kernel the projection is fixed; on the kernel complement the
spectral gap from (22.19) makes the two inverse-resolvent blocks tend to
zero and the off-diagonal block tend to zero in norm. The even kernel
has rank one and the odd kernel is zero. Thus the relative graph
projection has K-class \(+1\) in the compact operators. Therefore the KK
morphism \[
\delta_n=[e_0]^{-1}[e_1]\,m_{L^2\mathbb R^n}
\in KK(C_0(\mathbb R^{2n}),\mathbb C)
\tag{22.25}
\] satisfies \(\beta_n\delta_n=1_{\mathbb C}\).

Lemma 22.7 and actual semisplit KK exactness justify the inverse in
(22.25). The outward Bott element already has a proved two-sided KK
inverse \(\eta_n\), from the Bott/Clifford Morita proofs, with precisely
the convention fixed in the preceding index lesson's Euclidean
calculation. Consequently
\(\delta_n=(\eta_n\beta_n)\delta_n=\eta_n(\beta_n\delta_n)=\eta_n\).
This is a bivariant equality; it is not a classification of E-classes by
one integer. Lemma 22.5 gives
\(\operatorname{Ind}^E_{\mathbb R^n}=\mathfrak c(\delta_n)\). Applying
the product-preserving functor to the actual two KK inverse identities
proves both E inverse identities. For \(n=0\) this is the identity class
of a point. \(\square\)

The usual Weyl quantization deformation of \(C_0(\mathbb R^{2n})\) is
another presentation of this Euclidean index class. Its presentation by
Weyl kernels is stated here; (22.16)--(22.25) prove the index and Bott
normalization using the displayed positive and heat kernels.

\subsection{5. Embeddings and normal
deformation}\label{embeddings-and-normal-deformation}

For a smooth embedding \(i:M\hookrightarrow N\), the deformation
description uses two groupoids. The first is the vector translation
groupoid \[
J_i=i^*TN\rtimes_{di}TM,\qquad
r(x,\zeta,v)=(x,\zeta),\quad
s(x,\zeta,v)=(x,\zeta+di_xv).
\tag{22.26}
\] Scaling the action by \(\tau\) gives zero algebra
\(C_0(T^*M\oplus i^*TN)\) and positive algebra \(C^*(J_i)\). The second
has positive arrows \(N\times M\times M\), and zero arrows \(J_i\). In
local coordinates, its arrows approaching zero are \[
(X,\zeta,V,\tau)\longmapsto
(i(X)+\tau\zeta,X,X-\tau V,\tau).
\tag{22.27}
\] The source's normal coordinate tends to \(\zeta+di_XV\), since \[
\frac{i(X)+\tau\zeta-i(X-\tau V)}{\tau}
=\zeta+\int_0^1di_{X-s\tau V}V\,ds.
\tag{22.28}
\] Thus the sign and source in (22.26) are fixed by the actual boundary
chart, not by a choice of abstract normal-bundle notation. At positive
parameter the pair factor gives \(C_0(N)\otimes\mathcal K(L^2M)\). The
corresponding Haar factor is \(\tau^{-\dim M}\).

The chart, Haar, fibre and continuity proofs for these two deformations
are the explicit linear and differential-groupoid proofs in
\href{../KT-CP/KT-CP-17.html}{\emph{The tangent groupoid and deformation
to the normal cone}, ``Deformation classes and wrong-way maps''}. Apply
Theorem 22.2 to their cone extensions, rather than inferring a KK
equivalence for an arbitrary extension. This gives \[
\theta_i\in E(C_0(T^*M\oplus i^*TN),C^*(J_i)),\qquad
\gamma_i\in E(C^*(J_i),C_0(N)\otimes\mathcal K).
\tag{22.29}
\] A full K-orientation, including its determinant line and grading,
supplies the Thom factor on \(T^*M\oplus i^*TN\) in the convention of
\href{KT-KK-16.html}{\emph{Wrong-way maps for K-oriented maps}}.
Combining that factor, (22.29) and the stable pair-algebra corner gives
the deformation formulation of the oriented embedding class.

We state the comparison with the wrong-way class of that lesson: the
resulting E-class is \(\mathfrak c(i!)\), with the same full
orientation. The proof of this deformation comparison is not given here
and it is not an input to (22.5), (22.14), (22.15) or (22.25). Its free
primary locator is Connes's author edition, Chapter II, Section 6,
Propositions 1 and 3 and Theorem 7. The manifold construction,
composition and projection formulas themselves are the complete proofs
of the preceding wrong-way lesson.

The E-theoretic assembly map is a further use of deformation classes.
Its general construction and theorems for groups or groupoids are named
here only; none is used to establish the analytic index above.

\subsection{6. Exercises and complete
solutions}\label{exercises-and-complete-solutions}

\textbf{Exercise 1.} For a point-norm continuous path \(h_s:A\to B\) of
homomorphisms, construct its deformation extension and compute the
E-class. Determine when the most direct construction is a continuous
field.

\textbf{Solution.} Put \[
D_h=\{(a,b)\in A\oplus C([0,1],B):b(0)=h_0(a)\}.
\tag{22.30}
\] Projection onto \(a\) is onto, with kernel \(C_0((0,1],B)\). The
continuous lift \((a,s\mapsto h_s(a))\) makes (22.5) equal to
\(h_{1/t}\). This is asymptotically equivalent to \(h_0\), because
\(h_s(a)\to h_0(a)\) at every \(a\). The given ordinary homomorphism
path, suspended and stabilized, identifies \(h_0\) with \(h_1\). Hence
\(d_{D_h}\) is the endpoint homomorphism class.

For this direct field, the positive fibre norm of \((a,b)\) tends to
\(\|h_0(a)\|\), whereas the zero fibre norm is \(\|a\|\). It is a
continuous field precisely when \(h_0\) is isometric. An injective
C\emph{-homomorphism is isometric by the faithful C}-representation norm
theorem. If \(h_0\) has a nonzero kernel, \((a,0)\) for such an \(a\)
exhibits a norm jump. In particular a constant zero path on
\(A=\mathbb C\) is a counterexample to an unrestricted continuous-field
assertion. The deformation-extension construction and its E-class above
still apply. This distinction supplies the hypothesis required by the
continuous-field version of the exercise.

\textbf{Exercise 2.} Construct the tangent-groupoid asymptotic morphism
using a continuous section of zero evaluation.

\textbf{Solution.} Apply the countable weighted lift and residual series
(22.3)--(22.4) to \(e_0:D_M\to C_0(T^*M)\). For \(a\), evaluate its
lifted section at \(\tau=1/t\) in the pair-kernel representation on
\(L^2M\). Its multiplication error is exactly \[
e_{1/t}\bigl(\sigma(ab)-\sigma(a)\sigma(b)\bigr),
\] whose interior section is in \(C_0((0,1],\mathcal K(L^2M))\). Its
norm tends to zero; linearity and adjoint errors have the identical
proof. The suspension is the pointwise section \(u\mapsto\sigma(f(u))\),
whose errors converge uniformly by compact range and finite norm nets.
Compose with a rank-one corner. Another continuous section differs in
the same cone ideal and gives the same class. This is a family on every
symbol, not just on a smooth dense subalgebra.

\textbf{Exercise 3.} Compute the tangent deformation of a point.

\textbf{Solution.} There is one arrow at every parameter, all structure
maps are identity, and its Haar measure is one. The algebra is
\(C([0,1])\), with zero evaluation and positive scalar evaluation; its
ideal is \(C_0((0,1])\). Choose the constant lift of a scalar. The
family (22.5) is the identity of \(\mathbb C\). Its suspension and
stable corner are the identity E-class. Its K-map is \(1\mapsto1\) on
\(\mathbb Z\), with the even rank-one convention.

\textbf{Exercise 4.} Identify the Euclidean E-index as the inverse of
the outward Bott element.

\textbf{Solution.} Lemma 22.7 makes the zero evaluation a
KK-equivalence. Lemma 22.8 and (22.23)--(22.24) lift the outward Bott
graph projection. Its positive evaluation has exactly one even Gaussian,
no odd kernel and a gap on the complement, so the rescaling graph
homotopy gives \(\beta_n\delta_n=1\). The actual two-sided KK Bott
inverse \(\eta_n\) forces \(\delta_n=\eta_n\) by inverse uniqueness.
Formula (22.10) and Lemma 22.5 identify the direct E-deformation class
with \(\mathfrak c(\delta_n)\). Thus both its products with
\(\mathfrak c(\beta_n)\) are the corresponding identity E-classes.
Exterior grading and the complex coordinate \(x+i\eta\) fix the positive
sign; using the ordered right-odd coordinate convention would require
the simultaneous sign conversion of the element and its inverse.

\subsection{What this lesson does not
prove}\label{what-this-lesson-does-not-prove}

The general deformation comparison for manifold wrong-way maps is stated
in Section 5. The Weyl presentation and E-theoretic assembly are further
viewpoints, not proof premises. The analytic index of an uncontrolled
elliptic operator on a noncompact manifold is not asserted. Theorems
22.2, 22.4, 22.6 and 22.9 establish the actual deformation class, its
K-action, the closed operator index and the Euclidean bivariant inverse.

\subsection{Free mathematical reading}\label{free-mathematical-reading}

A. Connes,
\href{https://alainconnes.org/wp-content/uploads/book94bigpdf.pdf}{\emph{Noncommutative
Geometry}, freely posted author edition}, Chapter II, Section 5 and
Appendix B, for the tangent-groupoid and asymptotic deformation
construction; Section 6 for the normal-deformation viewpoint.

Y. Li,
\href{https://liyuezhao.github.io/notes/groupoid_notes.pdf}{\emph{Groupoid
C}-algebras*, open Leiden seminar notes, 7 February 2024}, Section 11,
for tangent groupoids and the two index maps.

B. Blackadar,
\href{https://www.bruceblackadar.com/Mathematics/book6.pdf}{\emph{K-Theory
for Operator Algebras}, author-posted corrected second edition}, Section
25.5, for extension classes in E-theory. The direct proofs above
distinguish arbitrary deformation extensions from semisplit KK
comparisons.

\end{document}
